\section{Backend - Indexing}

\subsection{Panoramica}

Il modulo di indicizzazione esplora il contenuto del pacchetto DiP e orchestra un flusso dati che va dal \codeword{PackageReaderPort} ai repository, garantendo persistenza efficiente delle entità estratte.
L'orchestrazione è gestita da \codeword{IndexDipUC}, che riceve in modalità asincrona tramite \codeword{AsyncGenerator} le entità del pacchetto (DiP, DocumentClass, Process, Document, File) dal \codeword{PackageReaderPort}. Ogni entità viene elaborata e persistita nel corrispettivo repository (\codeword{dipRepository}, \codeword{documentClassRepository}, \codeword{processRepository}, \codeword{documentRepository}, \codeword{fileRepository}).
Per i documenti, è previsto un trattamento speciale: il file principale (non gli allegati) subisce un processo di word embedding. L'estrazione del testo e la generazione dei vettori sono orchestrate da \codeword{EmbeddingService}, che interagisce con \codeword{WordEmbeddingPort} per convertire il contenuto testuale in vettori di embedding. Questi vettori vengono successivamente salvati in \codeword{VectorRepository}.

L'indicizzazione avviene in modalità asincrona, garantendo un utilizzo efficiente della memoria: le entità vengono elaborate e salvate man mano che vengono estratte dal pacchetto, senza attendere il completamento dell'esplorazione dell'intero contenuto.

Il modulo espone una inbound port per l'avvio del processo di indicizzazione, attualmente utilizzata da \codeword{ApplicationBootstrap} in \codeword{main.ts}. Poiché il software è distribuito all'interno del pacchetto DiP stesso, l'indicizzazione viene eseguita automaticamente durante il bootstrap dell'applicazione, senza richiedere intervento manuale dell'utente.

\subsection{Componenti}

\newpage
\begin{figure}[H]
    \centering
    \includegraphics[width=0.8\textwidth]{../../../Docs/out/3_PB/DocumentiEsterni/DiagrammiSpecificaTecnica/generated_/architecture/c1-c3/C3-main-indexing-flow-no-repos/C3_Main_Indexing_Flow_No_Repos.png}
        \caption{Diagramma C3 - Componenti del modulo di indicizzazione}
    \label{fig:C3_componenti_no_repos}
\end{figure}

\begin{figure}[H]
    \centering
    \includegraphics[width=1\textwidth]{../../../Docs/out/3_PB/DocumentiEsterni/DiagrammiSpecificaTecnica/generated_/architecture/c1-c3/C3-main-indexing-flow-repos/C3_Main_Indexing_Flow_Repos.png}
        \caption{Diagramma C3 - Componenti del modulo di indicizzazione}
    \label{fig:C3_componenti_repos}
\end{figure}
\newpage
\subsubsection{ApplicationBootstrap, IndexDipUC}

\begin{figure}[H]
    \centering
    \includegraphics[width=0.6\textwidth]{../../../Docs/out/3_PB/DocumentiEsterni/DiagrammiSpecificaTecnica/C4/C4_Orchestration/Class Diagram - Main Indexing Flow - Orchestration.png}
        \caption{Diagramma C4 - Dettaglio dell'orchestrazione del processo di indicizzazione}
    \label{fig:C4_orchestrazione_embedding}
\end{figure}

Prima di tutto è necessario capire la relazione tra le varie entità presenti all'interno di un pacchetto DiP:
\begin{itemize}
    \item \textbf{DiP}: è il pacchetto che viene indicizzato, contiene al suo interno tutte le altre entità.
    \item \textbf{DocumentClass}: rappresenta un aggregato di documenti, ma di fatto nel contesto di un pacchetto DiP rappresenta un aggregato di processi. Essendo il documento gestito all'interno di un processo, è stato aggiunto questo livello di innestamento per rappresentare questa relazione. Un DiP può contenere più DocumentClass.
    \item \textbf{Process}: rappresenta un processo, che contiene al suo interno più documenti.
    \item \textbf{Document}: rappresenta un documento, che contiene al suo interno più file.
    \item \textbf{File}: entità atomica contenuta all'interno di un documento. Può essere il documento principale del file (di cui ne esiste sempre uno e solo uno), oppure un allegato.
\end{itemize}

Viene spiegato in seguito come queste entità vengono gestite dal modulo di indicizzazione.

\paragraph{ApplicationBootstrap}:
\begin{itemize}
    \item \textbf{Descrizione}: adapter responsabile dell'avvio del processo di indicizzazione, utilizzato in \codeword{main.ts} per gestire il lifecycle dell'applicazione Electron.
    \item \codeword{\textbf{+ bootstrap()}: Promise<void>}: metodo che avvia il processo di indicizzazione, invocando \codeword{IndexDipUC.execute()}.
\end{itemize}

\paragraph{IndexDipUC}:
\begin{itemize}
    \item \textbf{Descrizione}: Use Case che orchestra il processo di indicizzazione, coordinando i flussi di dati in arrivo da \codeword{PackageReaderPort} verso i repository, effettuando l'embedding dei documenti tramite \codeword{WordEmbeddingPort}
    \item \codeword{\textbf{+execute(dipPath: string)}: Promise<IndexResult>}: metodo che avvia il processo di indicizzazione, accettando in ingresso il percorso del pacchetto DiP da indicizzare. Restituisce un oggetto \codeword{IndexResult} che rappresenta l'esito dell'operazione.
    \item \codeword{\textbf{-indexDip()}: Promise<IndexResult>}: Riceve l'entità DiP da \\
    \codeword{PackageReaderPort.readDip()} e la persiste nel corrispettivo repository
    \codeword{dipRepository}.
    \item \codeword{\textbf{-indexDocumentClasses()}: Promise<IndexResult>}: Riceve le DocumentClass da \\\codeword{PackageReaderPort.readDocumentClasses()} tramite AsyncGenerator e le persiste nel repository \codeword{documentClassRepository}.
    \item \codeword{\textbf{-indexProcesses()}: Promise<IndexResult>}: Riceve i Processi da \\\codeword{PackageReaderPort.readProcesses()} tramite AsyncGenerator e li persiste nel repository \codeword{processRepository}.
    \item \codeword{\textbf{-indexDocuments()}: Promise<IndexResult>}: Riceve i Documenti da \\\codeword{PackageReaderPort.readDocuments()} tramite AsyncGenerator e li persiste nel repository \codeword{documentRepository}. Per ogni documento, identifica il file principale, ne estrae il testo e lo converte in vettore tramite \codeword{WordEmbeddingPort}.
    \item \codeword{\textbf{-indexFiles()}: Promise<IndexResult>}: Riceve i File da\\
    \codeword{PackageReaderPort.readFiles()} tramite AsyncGenerator e li persiste nel\\
    repository \codeword{fileRepository}.
    \item \codeword{\textbf{-indexMainFileVector(file: File)}: Promise<IndexResult>}:\\
    Estrae il testo del file principale, lo converte in vettore tramite\\
    \codeword{WordEmbeddingPort} e lo salva in \codeword{VectorRepository}.
\end{itemize}

\paragraph{IndexResult}:
\begin{itemize}
    \item Descrizione: rappresenta l'esito del processo di indicizzazione, contenente una proprietà booleana \codeword{success} che indica se l'operazione è stata completata con successo o se si è verificato un errore.
    \item \codeword{\textbf{- success: boolean}}: indica l'esito del processo di indicizzazione, con \codeword{true} se l'operazione è stata completata con successo, o \codeword{false} se si è verificato un errore.
\end{itemize}

\subsubsection{EmbeddingService, WordEmbeddingPort}
Queste due componenti collaborano per gestire l'estrazione del testo dai file e la generazione dei vettori di embedding, che vengono poi restituiti a IndexDipUC per il salvataggio nella rispettiva \codeword{VectorRepository}.
\begin{figure}[H]
    \centering
    \includegraphics[width=0.8\textwidth]{../../../Docs/out/3_PB/DocumentiEsterni/DiagrammiSpecificaTecnica/C4/C4_Embedding/Class Diagram - Vertical Flow.png}
        \caption{Diagramma C4 - Dettaglio dell'orchestrazione del processo di indicizzazione}
    \label{fig:C4_orchestrazione_embedding}
\end{figure}

\paragraph{EmbeddingService}:
\begin{itemize}
    \item \textbf{Descrizione}: business service che si occupa di gestire la generazione di vettori a partire dal path del file principale di un documento. Questo servizio non si occupa della persistenza in quanto violerebbe il principio di Single Reponsibility.
    \item \codeword{\textbf{+ generateDocumentEmbedding(filePath: string)}: Promise<Float32Array | null>}: tramite \codeword{readTextFromStream} legge il contenuto del file principale del documento, che viene successivamente ridotto in chunk tramite \codeword{chunkText} e convertito in vettore tramite \codeword{WordEmbeddingPort}. Restituisce un array di float a 32 bit che rappresenta il vettore del documento, o null in caso di errore. Questo vettore è costruito come media dei vettori di ogni chunk del testo tramite \codeword{meanEmbedding}.
    \item \codeword{\textbf{+ setEmbeddingConfiguration(config: IEmbeddingConfiguration)}: void}: permette di aggiornare la configurazione per la generazione degli embedding, accettando in ingresso un oggetto che rispetta l'interfaccia \codeword{IEmbeddingConfiguration}.
    \item \codeword{\textbf{- readTextFromStream(filePath: string)}: Promise<string>}: legge il contenuto del file principale del documento a partire dal suo path tramite \\ \codeword{PackageReaderPort.readFileBytes}, restituendo una stringa di testo.
    \item \codeword{\textbf{- chunkText(text: string)}: string[]}: suddivide il testo in chunk di dimensione massima definita da \codeword{embeddingConfig.chunkSize}, restituendo un array di stringhe.
    \item \codeword{\textbf{- meanEmbedding(embeddings: Float32Array[])}: Float32Array | null}: calcola la media dei vettori di ogni chunk del testo, restituendo un array di float a 32 bit che rappresenta il vettore medio del documento, o null in caso di errore.
\end{itemize}

\paragraph{IEmbeddingConfiguration}:
\begin{itemize}
    \item \textbf{Descrizione}: interfaccia che definisce la configurazione per la generazione degli embedding.
    \item \codeword{chunkSize}: dimensione massima di ogni chunk di testo da convertire in vettore.
    \item \codeword{chunkOverlap}: sovrapposizione in byte tra un chunk e l'altro, per garantire continuità tra i chunk.
    \item \codeword{maxEmbeddingChunks}: numero massimo di chunk da convertire in vettore per ogni documento, permette di mantenere performance stabili con pacchetti contenenti file di dimensioni elevate.
\end{itemize}

\paragraph{WordEmbeddingPort}:
\begin{itemize}
    \item \textbf{Descrizione}: outbound adapter che permette di interagire con ONNX runtime per la generazione di vettori.
    \item \codeword{\textbf{+ initialize()}: Promise<void>}: inizializza il modello di embedding (paraphrase-multilingual-MiniLM-L12-v2) tramite la libreria \codeword{@xenova/transformers}. Questa operazione è asincrona e deve essere completata prima di poter generare embedding.
    \item \codeword{\textbf{+ generateEmbedding(texts: string[])}: Promise<Float32Array[]>}: genera i vettori di embedding per un array di stringhe di testo. Restituisce un array di Float32Array, uno per ogni testo in ingresso.
    \item \codeword{\textbf{+ isInitialized(): boolean}}: restituisce \codeword{true} se il modello è stato inizializzato, \codeword{false} altrimenti.
\end{itemize}

\subsubsection{LocalPackageReaderAdapter, DataMapper, DipParser, FileSystemProvider}

\begin{figure}[H]
    \centering
    \includegraphics[width=1\textwidth]{../../../Docs/out/3_PB/DocumentiEsterni/DiagrammiSpecificaTecnica/C4/C4_Main_Indexing_PackageReader/Class Diagram - Main Indexing Flow - Group 2.png}
        \caption{Diagramma C4 - Dettaglio dell'orchestrazione del processo di indicizzazione}
    \label{fig:C4_orchestrazione_reading}
\end{figure}

\paragraph{LocalPackageReaderAdapter}:
\begin{itemize}
    \item \textbf{Descrizione}: adapter che implementa \codeword{IPackageReaderPort}, responsabile della lettura del pacchetto DiP dal filesystem locale. Orchestr a il flusso di estrazione delle entità coordinando \codeword{DipParser}, \codeword{FileSystemProvider} e \codeword{DataMapper}.
    \item \codeword{\textbf{+ readDip(dipPath: string)}: Promise<Dip>}: legge il file \codeword{DiPIndex} dal percorso specificato, lo decodifica tramite \codeword{DipParser} e lo trasforma in entità \codeword{Dip} tramite \codeword{DataMapper}.
    \item \codeword{\textbf{+ readDocumentClasses(dipPath: string)}: AsyncGenerator<DocumentClass>}: genera in modo asincrono le entità \codeword{DocumentClass} estratte dal pacchetto. Per ogni classe documentale, legge il file XML corrispondente, lo decodifica e lo trasforma tramite \codeword{DataMapper}.
    \item \codeword{\textbf{+ readProcesses(dipPath: string)}: AsyncGenerator<Process>}: genera in modo asincrono le entità \codeword{Process} estratte dal pacchetto, leggendo e decodificando i file XML tramite \codeword{DipParser} e trasformandoli tramite \codeword{DataMapper}.
    \item \codeword{\textbf{+ readDocuments(dipPath: string)}: AsyncGenerator<Document>}: genera in modo asincrono le entità \codeword{Document} estratte dal pacchetto, gestendo la lettura dei metadati e la mappatura dei file allegati.
    \item \codeword{\textbf{+ readFiles(dipPath: string)}: AsyncGenerator<File>}: genera in modo asincrono le entità \codeword{File} estratte dal pacchetto, distinguendo tra file principale e allegati.
    \item \codeword{\textbf{+ readFileBytes(filePath: string)}: Promise<Buffer>}: legge il contenuto binario di un file dal filesystem, utilizzato da \codeword{EmbeddingService} per l'estrazione del testo.
\end{itemize}

\begin{figure}[H]
    \centering
    \includegraphics[width=1\textwidth]{../../../Docs/out/3_PB/DocumentiEsterni/DiagrammiSpecificaTecnica/C4/C4_Main_Indexing_Parser/Class Diagram - Main Indexing Flow - Group 2.png}
        \caption{Diagramma C4 - Dettaglio dell'orchestrazione del processo di indicizzazione}
    \label{fig:C4_orchestrazione_reading}
\end{figure}

\paragraph{IDipParser}:
\begin{itemize}
    \item \textbf{Descrizione}: interfaccia che permette l'attuazione dello \textbf{strategy pattern} per il parsing dei file XML. Il contratto è definito in modo da prevedere la massima flessibilità e facilità di implementazione dei diversi mapper. Viene accettato il contenuto del file (sia esso XML, JSON o altro) e restituito un oggetto JavaScript non tipizzato, che reifica il contenuto del file. L'attuale implementazione prevede il supporto di file XML.
    \item \codeword{\textbf{+ parse(rawContent: string)}: any}: effettua il parsing di una stringa XML e restituisce un oggetto JavaScript secondo lo schema fornito.
\end{itemize}

\paragraph{XmlDipParser}:
\begin{itemize}
    \item \textbf{Descrizione}: implementazione concreta di \codeword{IDipParser} che utilizza la libreria \codeword{fast-xml-parser} per la deserializzazione efficiente e type-safe dei file XML del pacchetto DiP.
    \item \textbf{Configurazione}: le opzioni di parsing sono configurate per gestire correttamente i namespace, gli attributi e le strutture complesse presenti negli schemi XSD forniti.
\end{itemize}

\begin{figure}[H]
    \centering
    \includegraphics[width=0.8\textwidth]{../../../Docs/out/3_PB/DocumentiEsterni/DiagrammiSpecificaTecnica/C4/C4_Main_Indexing_Flow_FsProvider/Class Diagram - Main Indexing Flow - Group 2.png}
        \caption{Diagramma C4 - Dettaglio dell'orchestrazione del processo di indicizzazione}
    \label{fig:C4_orchestrazione_reading}
\end{figure}

\paragraph{IFileSystemProvider}:
\begin{itemize}
    \item \textbf{Descrizione}: interfaccia che astrae l'accesso al filesystem, permettendo una lettura e una scrittura indipendente dalla piattaforma. Facilita inoltre il testing tramite mock.
    \item \codeword{\textbf{+ readFile(filePath: string)}: Promise<string>}: legge il contenuto testuale di un file dal percorso specificato, restituendo una stringa.
    \item \codeword{\textbf{+ readFileAsBuffer(filePath: string)}: Promise<Buffer>}: legge il contenuto binario di un file, restituendo un Buffer.
    \item \codeword{\textbf{+ listFiles(dirPath: string)}: Promise<string[]>}: elenca i file presenti in una directory, restituendo un array di percorsi relativi.
    \item \codeword{\textbf{+ exists(filePath: string)}: Promise<boolean>}: verifica l'esistenza di un file o directory al percorso specificato.
\end{itemize}

\paragraph{FileSystemProvider}:
\begin{itemize}
    \item \textbf{Descrizione}: implementazione concreta di \codeword{IFileSystemProvider} che utilizza le API native di Node.js (\codeword{fs} e \codeword{fs/promises}) per l'accesso al filesystem locale.
    \item \textbf{Caratteristiche}: supporta operazioni asincrone per evitare il blocco del thread principale, essenziale in un contesto Electron dove la responsività dell'UI è critica.
\end{itemize}

\begin{figure}[H]
    \centering
    \includegraphics[width=0.9\textwidth]{../../../Docs/out/3_PB/DocumentiEsterni/DiagrammiSpecificaTecnica/C4/C4_Main_Indexing_Flow_DataMapper/Class Diagram - Main Indexing Flow - Group 2.png}
        \caption{Diagramma C4 - Dettaglio dell'orchestrazione del processo di indicizzazione}
    \label{fig:C4_orchestrazione_reading}
\end{figure}

\paragraph{IDataMapper}:
\begin{itemize}
    \item \textbf{Descrizione}: interfaccia che definisce il contratto per la trasformazione dei dati grezzi estratti dai file XML nei modelli di dominio tipizzati. Utilizza il pattern \textbf{MapperRequest} per incapsulare i parametri di mapping e fornire contesto addizionale durante la trasformazione.
    \item \codeword{\textbf{+ map(request: MapperRequest<DipIndexXml>)}: Dip}: trasforma i dati grezzi del DiPIndex nelle entità \codeword{Dip}. Il \codeword{MapperRequest} fornisce il payload XML e il contesto di esecuzione.
    \item \codeword{\textbf{+ map(request: MapperRequest<DocumentClassXml>)}: DocumentClass}: trasforma i dati grezzi nelle entità \codeword{DocumentClass}, accedendo a metadati contestuali attraverso \codeword{MapperRequest}.
    \item \codeword{\textbf{+ map(request: MapperRequest<AiPXml>)}: Process}: trasforma i dati grezzi nelle entità \codeword{Process}, utilizzando il contesto fornito da \codeword{MapperRequest} per correlazioni tra entità.
    \item \codeword{\textbf{+ map(request: MapperRequest<DocumentXml>)}: Document}: trasforma i dati grezzi nelle entità \codeword{Document}, gestendo i metadati e la relazione con i file attraverso il contesto di \codeword{MapperRequest}.
    \item \codeword{\textbf{+ map(request: MapperRequest<FileXml>)}: File}: trasforma i dati grezzi nelle entità \codeword{File}, distinguendo tra file principale e allegati tramite le informazioni contestuali di \codeword{MapperRequest}.
\end{itemize}

\paragraph{MapperRequest<T>}:
\begin{itemize}
    \item \textbf{Descrizione}: interfaccia generica che incapsula i dati da trasformare e fornisce il contesto di esecuzione necessario per il mapping. Permette di passare informazioni addizionali al mapper senza modificare la firma dei metodi.
    \item \codeword{\textbf{+ payload: T}}: il dato grezzo estratto dal file XML da trasformare nel modello di dominio.
    \item \codeword{\textbf{+ context?: Record<string, any>}}: oggetto opzionale che contiene metadati contestuali (es. dipId, parentPath, processId) utili durante la trasformazione per stabilire correlazioni tra entità e arricchire il mapping con informazioni di relazione.
    \item \textbf{Vantaggi}: consente di mantenere le firme dei metodi stabili pur permettendo l'espansione del contesto senza breaking changes, facilita il testing tramite iniection del contesto, e rende esplicite le dipendenze contestuali del mapping.
\end{itemize}

\paragraph{DataMapper}:
\begin{itemize}
    \item \textbf{Descrizione}: implementazione concreta di \codeword{IDataMapper} che effettua la trasformazione dei dati grezzi nei modelli di dominio, applicando le regole di business specifiche del contesto DiP. Collabora con \codeword{DocumentMetadataHashMapper} per la normalizzazione dei metadati.
    \item \textbf{Responsabilità}:
    \begin{itemize}
        \item Mappatura delle strutture dati XML nelle entità di dominio fortemente tipizzate.
        \item Gestione delle relazioni tra entità (es. documento-file, processo-documentclass).
        \item Normalizzazione e validazione dei metadati estratti dal pacchetto.
        \item Conversione dei tipi primitivi (stringhe, date, etc.) nei corretti tipi TypeScript.
    \end{itemize}
\end{itemize}

\paragraph{DocumentMetadataHashMapper}:
\begin{itemize}
    \item \textbf{Descrizione}: utility class specializzata nella normalizzazione e nella trasformazione dei metadati dei documenti estratti dal pacchetto DiP. Converte la rappresentazione key-value irregolare in strutture coerenti per la persistenza.
    \item \codeword{\textbf{+ normalizeMetadata(rawMetadata: Record<string, any>)}: Metadata[]}: normalizza i metadati grezzi, gestendo variabilità nella struttura e nei formati, e restituisce un array di entità \codeword{Metadata} fortemente tipizzate.
    \item \textbf{Logica}: implementa regole di trasformazione specifiche per i metadati del contesto archivistico digitale, garantendo coerenza e correttezza dell'informazione.
\end{itemize}
